% ============================================================
% Appendix: Computational Details
% XCEPT Article 1 — All three sections
% ============================================================

\section{Computational Details}
\label{app:computational}

This appendix documents all modelling choices, parameter values, and
algorithmic decisions for the three analytical sections of the paper.
The full source code is available in the accompanying repository;
the three scripts \texttt{script\_section3.py}, \texttt{script\_section4.py},
and \texttt{script\_section5.py} reproduce every figure directly.

% ============================================================
\subsection{Section~3 — Band-Tribe Conflict Dynamics}
\label{app:sec3}

\subsubsection*{Reaction Network}

The model consists of five species
$\mathcal{S} = \{C, X, \textit{Res}, C\_\textit{Res}, G\}$
and eleven reactions $r_0$--$r_{10}$, loaded from
\texttt{data/Basic\_example.txt}.  Table~\ref{tab:s3_reactions} lists
each reaction, its biological interpretation, and its kinetic law.

\begin{table}[h]
\centering
\small
\caption{Band-tribe model reactions.  MM = Michaelis-Menten saturation;
         MA = mass-action.  $\alpha_S$ and $\alpha_R$ are scenario
         modifiers for solidarity and reactivity respectively.}
\label{tab:s3_reactions}
\begin{tabular}{lllll}
\toprule
ID   & Stoichiometry                        & Interpretation                 & Kinetic law         & Type \\
\midrule
$r_0$ & $\varnothing \to \textit{Res}$       & Constant resource inflow        & $\kappa_0$          & MA   \\
$r_1$ & $C + \textit{Res} \to C$             & Band consumption                & $\kappa_1 \cdot C \cdot \frac{\textit{Res}}{\textit{Res}+K_{r1}}$ & MM \\
$r_2$ & $C + \textit{Res} + C\_\textit{Res} \to C + 2\,C\_\textit{Res}$ & Solidarity: surplus growth & $\alpha_S\,\kappa_2 \cdot C \cdot \frac{\textit{Res}}{\textit{Res}+K_{r2}} \cdot C\_\textit{Res}$ & MM \\
$r_3$ & $C\_\textit{Res} \to \textit{Res}$   & Resilience redistribution       & $\kappa_3 \cdot \frac{C\_\textit{Res}}{C\_\textit{Res}+K_{r3}}$ & MM \\
$r_4$ & $C \to X$                             & Spontaneous displacement        & $\kappa_4 \cdot C$  & MA   \\
$r_5$ & $X + C\_\textit{Res} \to C$           & Solidarity: reintegration       & $\alpha_S\,\kappa_5 \cdot X \cdot C\_\textit{Res}$ & MA \\
$r_6$ & $2\,\textit{Res} \to \varnothing$     & Resource decay (second-order)   & $\kappa_6 \cdot \textit{Res}^2$ & MA \\
$r_7$ & $2\,C\_\textit{Res} \to \varnothing$  & Resilience decay (second-order) & $\kappa_7 \cdot C\_\textit{Res}^2$ & MA \\
$r_8$ & $X+C+\textit{Res} \to X+C+G$         & Reactivity: grievance generation & $\alpha_R\,\kappa_8 \cdot X \cdot C \cdot \frac{\textit{Res}}{\textit{Res}+K_{r8}}$ & MM \\
$r_9$ & $2G + C \to X$                        & Reactivity: displacement by grievance & $\alpha_R\,\kappa_9 \cdot G^2 \cdot C$ & MA \\
$r_{10}$ & $2G \to \varnothing$              & Grievance decay (second-order)  & $\kappa_{10} \cdot G^2$ & MA \\
\bottomrule
\end{tabular}
\end{table}

\subsubsection*{ODE Kinetic Parameters}

\begin{table}[h]
\centering
\caption{Base rate constants and Michaelis-Menten saturation constants for the ODE.}
\label{tab:s3_ode_params}
\begin{tabular}{lll}
\toprule
Parameter         & Value   & Description \\
\midrule
$\kappa_0$        & 0.200   & Constant resource inflow \\
$\kappa_1$        & 0.020   & Band consumption rate \\
$\kappa_2$        & 0.100   & Surplus growth (scaled by $\alpha_S$) \\
$\kappa_3$        & 0.010   & Resilience redistribution \\
$\kappa_4$        & 0.010   & Spontaneous displacement \\
$\kappa_5$        & 0.100   & Reintegration (scaled by $\alpha_S$) \\
$\kappa_6$        & 0.005   & Res decay (second-order) \\
$\kappa_7$        & 0.010   & $C\_\textit{Res}$ decay (second-order) \\
$\kappa_8$        & 0.010   & Grievance generation (scaled by $\alpha_R$) \\
$\kappa_9$        & 0.050   & Displacement by grievance (scaled by $\alpha_R$) \\
$\kappa_{10}$     & 0.005   & Grievance decay (second-order) \\
\midrule
$K_{r1}$          & 3.0     & MM saturation in $r_1$ (Res) \\
$K_{r2}$          & 3.0     & MM saturation in $r_2$ (Res) \\
$K_{r3}$          & 2.0     & MM saturation in $r_3$ ($C\_\textit{Res}$) \\
$K_{r8}$          & 3.0     & MM saturation in $r_8$ (Res) \\
\bottomrule
\end{tabular}
\end{table}

\subsubsection*{Two Scenarios}

The solidarity modifier $\alpha_S$ scales reactions $r_2$ and $r_5$;
the reactivity modifier $\alpha_R$ scales reactions $r_8$ and $r_9$.
Two contrasting scenarios are compared:

\begin{itemize}
    \item \textbf{Scenario A} (high solidarity, low reactivity):
          $\alpha_S = 2.0$, $\alpha_R = 0.5$.
    \item \textbf{Scenario B} (low solidarity, high reactivity):
          $\alpha_S = 0.3$, $\alpha_R = 3.0$.
\end{itemize}

Both scenarios share the initial condition
$\mathbf{x}_0 = (C, X, \textit{Res}, C\_\textit{Res}, G) = (10, 3, 5, 2, 1)$.

\subsubsection*{ODE Integration}

The ODE is integrated using \texttt{scipy.integrate.solve\_ivp} with the
LSODA method, tolerances $\texttt{rtol}=10^{-8}$, $\texttt{atol}=10^{-10}$,
and evaluation at $N = 2000$ evenly spaced points on $[0, T_{\text{end}}]$
with $T_{\text{end}} = 400$.
All denominators are protected by a floor of $10^{-9}$ to avoid division by zero.

\subsubsection*{Stochastic Simulation (Poisson Tau-Leaping)}

The stochastic simulation uses Poisson tau-leaping with a fixed step size
$\Delta t = 1.0$ for $N_{\text{steps}} = 300$ steps.
The propensity of reaction $j$ at state $\mathbf{x}$ is
\begin{equation}
    \lambda_j = \kappa^{\text{stoch}}_j \cdot \pi_j(\mathbf{x}),
    \qquad
    \pi_j(\mathbf{x}) = \min_{s:\,S^-_{sj}>0} \frac{x_s}{S^-_{sj}},
    \label{eq:s3_propensity}
\end{equation}
where $S^-_{sj}$ is the input stoichiometry (reactant coefficient) and
$\pi_j$ is the enabling degree (bottleneck reactant).  At each step,
firing counts are drawn as $n_j \sim \mathrm{Poisson}(\lambda_j \Delta t)$,
capped at $\lfloor\pi_j\rfloor$, and then rescaled by an iterative contention
resolver (maximum 20 iterations) that prevents any species from going negative.

The stochastic rate constants $\kappa^{\text{stoch}}$ differ from the ODE
constants because the enabling-degree propensity underestimates multiplicative
product kinetics.  At initial conditions, the ratio of the ODE rate to the
enabling degree can reach 6$\times$ for three-reactant reactions such as $r_2$.
Both parameter sets are independent calibrations of the same qualitative model.

\begin{table}[h]
\centering
\caption{Stochastic rate constants $\kappa^{\text{stoch}}$.  Modifier
         reactions carry the same $\alpha$ factors as the ODE.}
\label{tab:s3_stoch_params}
\begin{tabular}{lll}
\toprule
Reaction  & $\kappa^{\text{stoch}}$ & Note \\
\midrule
$r_0$     & 2.000   & Source (no reactant; $\pi_j=1$) \\
$r_1$     & 0.025   & ODE/enabling ratio $\approx 1$ at $\mathbf{x}_0$ \\
$r_2$     & 0.500   & Solidarity; $\pi=\min(C,\textit{Res},C\_\textit{Res})=2$ at $\mathbf{x}_0$ \\
$r_3$     & 0.020   & \\
$r_4$     & 0.010   & Matches ODE \\
$r_5$     & 0.100   & Solidarity; $\pi=\min(X,C\_\textit{Res})=2$ at $\mathbf{x}_0$ \\
$r_6$     & 0.050   & Second-order; $\pi=\textit{Res}/2$ \\
$r_7$     & 0.020   & Second-order \\
$r_8$     & 0.050   & Reactivity; $\pi=\min(X,C,\textit{Res})=3$ \\
$r_9$     & 0.100   & Reactivity; $\pi=\min(G/2,C)=0.5$ \\
$r_{10}$  & 0.010   & Grievance decay \\
\bottomrule
\end{tabular}
\end{table}

% ============================================================
\subsection{Section~4 — Chiefdom Emergence and Strategy}
\label{app:sec4}

\subsubsection*{Reaction Network}

The chiefdom model extends the band-tribe network by adding a hierarchy
species $H$, its resource pool $H\_\textit{Res}$, a police force $P_H$,
and an external aid signal $F_H$ (catalytic; never consumed).  The full
system has nine species
\begin{equation*}
    \mathcal{S} = \{C,\;X,\;H,\;\textit{Res},\;C\_\textit{Res},\;
                    H\_\textit{Res},\;G,\;P_H,\;F_H\}
\end{equation*}
and twenty-two reactions $r_0$--$r_{21}$ loaded from
\texttt{data/Intrmediate\_example.txt}.
Reactions $r_0$--$r_{10}$ are the band-tribe baseline (parameterised at
Scenario~A: $\alpha_S=2$, $\alpha_R=0.5$ baked into $\kappa^{\text{stoch}}$).
Reactions $r_{11}$--$r_{21}$ govern the chiefdom.

\subsubsection*{Initial Conditions}

\begin{equation*}
    \mathbf{x}_0 = (C,\;X,\;H,\;\textit{Res},\;C\_\textit{Res},\;
                    H\_\textit{Res},\;G,\;P_H,\;F_H)
                 = (500,\;50,\;50,\;10,\;50,\;100,\;30,\;10,\;0).
\end{equation*}
The chiefdom starts as a small perturbation ($H=50$, $H\_\textit{Res}=100$)
on top of a harmonious band-tribe baseline.

\subsubsection*{Stochastic Simulation}

Simulation uses Poisson tau-leaping with $\Delta t = 1.0$ and
$N_{\text{steps}} = 5000$, with $N_{\text{seeds}} = 10$ independent
replicates per scenario.
Foreign aid $F_H$ is injected adaptively: $F_H = 100$ whenever
$x_X \geq 50$ or $x_G \geq 10$; otherwise $F_H = 0$ (aid decays
immediately when conditions are not met).

\subsubsection*{Chiefdom Allocation Strategies}

Eight chiefdom-controlled reactions are scaled by allocation multipliers
$A \in \mathbb{R}_{>0}^8$ (index order: $r_{12}, r_{13}, r_{15}, r_{16},
r_{17}, r_{18}, r_{19}, r_{20}$):

\begin{table}[h]
\centering
\caption{Chiefdom allocation multipliers for the two strategies.
         Bold entries indicate reactions that are central to the mode
         definition ($\times$\,weight in mode vector shown in parentheses).}
\label{tab:s4_allocation}
\begin{tabular}{llrrl}
\toprule
Reaction & Role & $A_\text{protect}$ & $A_\text{exploit}$ & Mode membership \\
\midrule
$r_{12}$: $C+H+\textit{Res}\to H\_\textit{Res}$   & Res extraction       & \textbf{8.0} & 3.0          & Protection $(\times 2)$; Exploitation $(\times 1)$ \\
$r_{13}$: $H+C\_\textit{Res}\to H\_\textit{Res}$ & $C\_\textit{Res}$ extraction & 1.0 & \textbf{10.0} & Exploitation $(\times 1)$ \\
$r_{15}$: $C+H+H\_\textit{Res}\to 2H$             & Hierarchy expansion  & 5.0          & \textbf{8.0} & Exploitation $(\times 2)$ \\
$r_{16}$: $H+X+H\_\textit{Res}\to P_H$            & Militarize displaced & \textbf{8.0} & 0.5          & Protection $(\times 1)$ \\
$r_{17}$: $P_H+H\_\textit{Res}\to P_H$            & Police upkeep        & 1.0          & 1.0          & (neutral) \\
$r_{18}$: $2H\_\textit{Res}\to H\_\textit{Res}+C\_\textit{Res}$ & Redistribution & \textbf{5.0} & 0.05 & Protection $(\times 1)$ \\
$r_{19}$: $P_H+G\to P_H$                          & Grievance suppression & 2.0         & 5.0          & (background) \\
$r_{20}$: $P_H\to C$                               & Demobilize police    & 1.0          & 0.1          & (background) \\
\bottomrule
\end{tabular}
\end{table}

\subsubsection*{Mode Definitions}

Five canonical process-vector modes are defined as directions in
reaction-rate space.  Each mode is a non-negative vector
$\mathbf{d} \in \mathbb{R}_{\geq 0}^{22}$; entries are reaction weights:

\begin{description}
    \item[Community recovery:]
          $\mathbf{d}_{\text{comm}} = r_0 + r_2 + r_5$. \\
          Self-reinforcing band-tribe loop: resource inflow $\to$ $C\_\textit{Res}$
          surplus $\to$ reintegration of displaced.

    \item[Conflict amplification:]
          $\mathbf{d}_{\text{conf}} = r_4 + 2r_8 + r_9$. \\
          Destabilising spiral: $C\to X$, grievance accumulation, and
          grievance-driven further displacement.  $r_8$ is double-weighted
          because it drives both grievance accumulation and feeds $r_9$.

    \item[Chiefdom protection:]
          $\mathbf{d}_{\text{prot}} = 2r_{12} + r_{16} + r_{18}$. \\
          Fuel path ($2r_{12}$): double-weight Res extraction builds $H\_\textit{Res}$;
          $r_{13}$ is \emph{not} in this mode (no $C\_\textit{Res}$ drain).
          Redistribution path ($r_{18}$, KEY): returns wealth to community.
          Protective-force path ($r_{16}$): absorbs displaced into police.

    \item[Chiefdom exploitation:]
          $\mathbf{d}_{\text{expl}} = r_{12} + r_{13} + 2r_{15}$. \\
          Dual extraction fuel ($r_{12}+r_{13}$) feeds hierarchy expansion
          cycle ($2r_{15}$): $C+H+H\_\textit{Res}\to 2H$.
          $r_{18}$ is suppressed; no wealth returns to the community.

    \item[Foreign aid:]
          $\mathbf{d}_{\text{aid}} = r_{21}$. \\
          Catalytic $H\_\textit{Res}$ injection.  Excluded from all projection
          computations so that the four main modes are comparable across
          aid and no-aid scenarios.
\end{description}

\subsubsection*{Mode Projection Analysis}

For a non-overlapping window of $w$ steps ending at step $t$, the
per-step average process vector is
\begin{equation}
    \bar{\mathbf{v}} = \frac{1}{w}\sum_{s=t-w+1}^{t} \mathbf{n}_s,
\end{equation}
where $\mathbf{n}_s \in \mathbb{Z}_{\geq 0}^{22}$ is the reaction-firing
count vector at step $s$.  For mode $\mathbf{d}$ with unit vector
$\hat{\mathbf{d}} = \mathbf{d}/\|\mathbf{d}\|$, two summary statistics are
computed:
\begin{equation}
    \text{proj} = \hat{\mathbf{d}} \cdot \bar{\mathbf{v}},
    \qquad
    \text{cos-sim} = \frac{\hat{\mathbf{d}} \cdot \bar{\mathbf{v}}}
                          {\|\bar{\mathbf{v}}\|}.
\end{equation}
The projection magnitude measures absolute mode activity; the cosine
similarity measures mode dominance relative to all other reactions firing
simultaneously.  The scatter plots in Section~4 show one point per
non-overlapping block per seed for each window size
$w \in \{10, 50, 200, 500\}$.

\subsubsection*{Timescale Mode Classification}

Each species group is classified into one of four qualitative modes
based on the net change of each species within the window:
\begin{itemize}
    \item \textbf{problem} — all $\Delta x_s < 0$ (net decline);
    \item \textbf{overproduction} — all $\Delta x_s > 0$ (net growth);
    \item \textbf{challenge} — mixed (some species grow, some decline);
    \item \textbf{steady-state} — all $|\Delta x_s| < \tau_{\text{tol}}$
          (near-zero change, $\tau_{\text{tol}} = 0.05$).
\end{itemize}
Three species groups are tracked:
$M_P = \{C, C\_\textit{Res}\}$ (community production),
$M_H = \{H, H\_\textit{Res}\}$ (chiefdom),
$M_C = \{X, G\}$ (conflict indicators).
Cell opacity in the raster figures encodes the primary species' share of
the total human population $C + H + X$, with parameters
$\alpha_{\min} = 0.25$, $\alpha_{\max} = 1.0$, and full opacity at
a share $\geq 0.5$.

\subsubsection*{H\textsubscript{Res} Community-Benefit Routing}

The routing fraction measures what fraction of total $H\_\textit{Res}$
outflow benefits the community:
\begin{equation}
    \rho = \frac{n_{16} \cdot 1 + n_{18} \cdot 2}
               {n_{11} \cdot 1 + n_{15} \cdot 1 + n_{16} \cdot 1
                + n_{17} \cdot 1 + n_{18} \cdot 2},
\end{equation}
where $n_j$ is the cumulative firing count of reaction $j$ in the window
and the integer weights are the $H\_\textit{Res}$ input stoichiometry
coefficients.  $\rho = 1$ means all $H\_\textit{Res}$ flows to redistribution
and peacekeeping; $\rho = 0$ means all flows to elite maintenance or
expansion.  Summary statistics use $w = 200$ averaged over the last 50\%
of simulation steps.

% ============================================================
\subsection{Section~5 — Structural Markov Graph and ESMO Analysis}
\label{app:sec5}

\subsubsection*{Overview}

Section~5 analyses transitions between organisational states using
Extreme Stochastic Mode Objects (ESMOs).  An ESMO is a vertex of the
polytope defined by all non-negative flow vectors
$\mathbf{v} \geq \mathbf{0}$ that are consistent with the stoichiometry
of a given transition signature.  The number of ESMOs for a source--target
state pair is used as a proxy for the structural accessibility of that
transition.

\subsubsection*{Organisational Lattice}

Four organisational types form a cover lattice under the inclusion relation:
\begin{equation*}
    \text{Tribe} \subset \text{Chief},\quad
    \text{Tribe} \subset \text{State},\quad
    \text{Chief} \subset \text{ChiefState},\quad
    \text{State} \subset \text{ChiefState}.
\end{equation*}
Each organisation can be in peace ($P$) or conflict ($C$) state, yielding
eight nodes: TrP, TrC, ChP, ChC, StP, StC, CSP, CSC.
Inter-organisational transitions follow the four cover edges
(\emph{up} = smaller $\to$ larger org; \emph{down} = larger $\to$ smaller),
and each directed cover edge carries two \emph{combo} types (PP/PC for peace
sources; CP/CC for conflict sources).

\subsubsection*{ESMO Computation}

For each (edge, direction, combo) triple, the LP solver enumerates all
vertices of the transition-feasibility cone.  Vertex enumeration is
performed using pyCOT's \texttt{Persistent\_Modules\_Generator}, which
builds the reaction network from the transition signature and calls an
external vertex-enumeration routine.  The output is the number of vertices
$N_{\text{ESMO}}$ — a count of structurally distinct pathways through which
the transition can be realised.

Results are cached in \texttt{outputs/complex\_COT\_ESMO/transition\_summary.csv}
and \texttt{outputs/intra\_org\_esmo/esmo\_cache\_intra.csv} to avoid
recomputation (which can take several hours).  The cache files are read
directly by \texttt{script\_section5.py} to regenerate all figures.

\subsubsection*{Conditional Relative Frequency ($f$-value)}

For a directed cover edge and a given source-state type, the conditional
relative frequency is:
\begin{equation}
    f(\text{combo}) = \frac{N_{\text{ESMO}}(\text{combo})}
                          {N_{\text{ESMO}}(\text{combo}) + N_{\text{ESMO}}(\text{complement})},
\end{equation}
where the complement swaps only the target-state component
(PP$\leftrightarrow$PC; CP$\leftrightarrow$CC).
Thus $f(\text{PP}) + f(\text{PC}) = 1$ for peace sources and
$f(\text{CP}) + f(\text{CC}) = 1$ for conflict sources on the same
directed edge.  Each node has two outgoing cover edges; summing $f$-values
across both edges gives 2, not 1.

\subsubsection*{Normalised Transition Probability ($P$-value)}

To obtain probabilities that sum to 1 per source node across all four
outgoing transitions, we compute:
\begin{equation}
    P(u \to v) = \frac{N_{\text{ESMO}}(u \to v)}
                      {\displaystyle\sum_{v':\,(u,v')\in E} N_{\text{ESMO}}(u \to v')},
\end{equation}
where the denominator sums over all four outgoing directed transitions
from node $u$.  The $P$-values reported in the Markov graph and in the
inter-organisational bar charts (Section~5, Figures~5--7) are all
computed with this normalisation.

\subsubsection*{Markov Graph Visualisation}

The Markov graph is rendered in two formats by \texttt{plot\_markov\_hierarchy.py}:
\begin{itemize}
    \item An interactive HTML file (\texttt{markov\_hierarchy.html}) using
          the PyVis library.
    \item A static PNG (\texttt{fig\_markov\_loops.png}) using Matplotlib.
\end{itemize}

Arrow properties are driven by the $P$-value via a single unified formula:
\begin{equation}
    \text{line width} = 12\,P,
    \qquad
    \text{opacity} = 2.2\,P.
\end{equation}
With $P \in [0.10, 0.40]$, this maps to line widths $[1.2, 4.8]$ and
opacities $[0.22, 0.88]$.  All 32 arrows use the same formula with curvature
$\text{rad} = 0.15$; bidirectional pairs curve to opposite sides automatically
because the arc direction is relative to the travel direction.

Arrow colours encode the combo type:
\begin{center}
\begin{tabular}{lll}
\toprule
Combo & Direction & Colour \\
\midrule
PP & peace $\to$ peace     & forest green \texttt{\#3D7A4A} \\
PC & peace $\to$ conflict  & amber \texttt{\#C8810A} \\
CP & conflict $\to$ peace  & steel blue \texttt{\#2563A6} \\
CC & conflict $\to$ conflict & brick red \texttt{\#A83232} \\
\bottomrule
\end{tabular}
\end{center}

The configuration variable \texttt{SHOW\_VALUE = 'P'} controls which
value appears as edge labels (options: \texttt{'P'}, \texttt{'f'}, or
\texttt{'both'}).  The variable \texttt{WEIGHT\_COL = 'P'} independently
controls which value drives arrow thickness and opacity.

\subsubsection*{Dominant Structural Cycle}

Greedy path-tracing from TrC (the most conflict-prone tribal node)
following the maximum-$P$ outgoing transition at each step reveals
the dominant 4-cycle:
\begin{equation}
    \text{TrC} \xrightarrow{P=0.344,\;\text{PC}} \text{ChC}
    \xrightarrow{P=0.374,\;\text{CC}} \text{CSC}
    \xrightarrow{P=0.352,\;\text{CC}} \text{StP}
    \xrightarrow{P=0.308,\;\text{CP}} \text{TrC}.
    \label{eq:dominant_cycle}
\end{equation}
This cycle captures a structural attractor in which tribal conflict escalates
through hierarchical layers before a fragile state-level peace formation
collapses back to tribal conflict.
